\section{Dimensionamiento de fibras y capacidad}
\label{sec:dimensionamiento}

\subsection{Fibras y cables}
\label{sec:fibras}
% Consigna 6.1: fibras necesarias en cada tramo para la cobertura final. Diferenciar:
% primera etapa, crecimiento, reserva y otras aplicaciones.
% Seleccionar capacidades comerciales de cable y presentar una tabla por tramo con
% longitud, cantidad de fibras y asignación.
\pendiente{tabla por tramo: longitud, fibras y asignación (uso, crecimiento, reserva).}

\subsection{Puertos GPON y ancho de banda}
\label{sec:puertos}
% Consigna 6.2: definir y justificar:
% - Máximo de abonados por puerto GPON.
% - Relación de división total y etapas de división.
% - Puertos necesarios en la primera etapa y en la cobertura final.
% - Capacidad compartida y demanda estimada en hora pico.
% - Criterio de simultaneidad o sobrecontratación.
% Distinguir la relación de división óptica del criterio de sobrecontratación y explicar
% cómo ambos condicionan la capacidad ofrecida.
\pendiente{abonados por puerto, relación de división, cantidad de puertos y demanda en hora pico.}
